
This chapter closes the dissertation: what each objective was met with and how the two research questions are answered (\S9.1), the contributions (\S9.2), each limitation paired with the experiment that would settle it (\S9.3), and a personal reflection on how the project was actually conducted (\S9.4).

\section{Summary of the dissertation}

The aim was to determine whether spatial-relationship annotation for robot-acquired images can be automated, and whether the resulting labels are good enough to replace the human ones. Six objectives carried that aim, and each is met with evidence rather than assertion.

Each is met with evidence rather than assertion, and the evidence is already reported: a pipeline annotating all 836 images in about five minutes on a 6 GB consumer GPU with no human deciding any label, writing the dataset's own formats byte-compatibly and verified by a zero-error load-write round trip (\textbf{O1}); written operational definitions of all seven predicates with every threshold fitted on annotator groups 0 to 5 alone, \texttt{near} generalising to the held-out annotator at recall 1.00 and the support rules reaching held-out F1 0.87 (\textbf{O2}); per-predicate fidelity against all 8,926 human triplets with three trivial baselines, a vision-language model under the same definitions (\S4.13), nine ablations, audited precision and cluster-bootstrap intervals, giving mean recall 0.85 and 0.76 on annotators no threshold ever saw (\textbf{O3}); every one of the 1,689 missed triplets attributed to a cause, leaving roughly 7\% attributable to genuine tool error (\S4.10, \S7.2) (\textbf{O4}); a controlled experiment isolating the label source, at 0.76 mean recall against held-out human gold for the automatic arm, 0.30 human and 0.36 self-trained (\textbf{O5}); and the same comparison repeated in the source paper's own framework with a frozen detector and three seeds per arm, carried one link further to a planner (\textbf{O6}). O6 is met in the sense that matters for an honest answer: the heavyweight test does not agree with the lightweight one, the disagreement is localised to the annotators whose labels Chapter 4 convicted of measured defects, and both readings are reported rather than one selected.

Section 1.2.2 fixed what would count as an answer before any result was reported, so both questions are settled against those criteria rather than a measure chosen afterwards. Section 7.1 argues the answers; what follows is the verdict and its conditions.

\textbf{RQ1 is answered yes on five of seven predicates and qualified on two.} The criterion was per-predicate recall on annotator groups no threshold ever saw, judged against two references fixed in advance, and both are met: the trivial baselines are beaten by a wide margin (0.85 against 0.14), and the tool's mean agreement with the consistent annotators, 0.869, falls inside the [0.74, 0.92] interval within which those annotators can be shown to agree with one another (\S4.6), so the tool is not distinguishable from a tenth annotator. The second condition, that labels beyond the human record survive audit, holds for the lateral, proximity and depth-decided predicates at audited precision near 1.0 and for support at roughly 0.9 after the contact rule. The qualification is \texttt{in front of} and \texttt{behind} at 0.64/0.66 pooled: not a failure of the criterion but a disagreement about the words, since \S4.12 shows the tool reproducing its own verdict across viewpoints 0.955 of the time while two annotator groups applied the opposite convention. A per-predicate answer was required precisely so this could not hide inside a mean.

\textbf{RQ2 is answered yes at two of the three levels the criterion named, and the third disagrees.} The controlled classifier gives 0.76 against 0.30 and the planner 22 of 25 against 0 of 25 with no relations, both in the harder direction the criterion specified: held-out \textit{human} gold, the rival source's own yardstick. The benchmark does not agree, the human arm ranking better on mR@100 in the source paper's own framework, 0.326 against 0.278 over three seeds. Section 1.2.2 committed to reporting that as disagreement rather than resolving it favourably, and Chapter 6 does: the human arm's lead sits entirely on the two test annotators carrying measured labelling defects and vanishes on the one without (0.308 against 0.307), while the auto arm covers sixty times more of the relation types the manual annotation never recorded. An unqualified yes required agreement across all three, so the honest verdict is conditional.

The condition is legible, and it is the sentence the evidence supports: \textbf{automatic labels are the better supervision wherever ground truth means geometry; human labels remain better wherever ground truth means human annotation practice.} Robot planning needs the former, which is why the planner result and the benchmark result point in opposite directions without contradicting each other.

\section{Research contributions}

\textbf{For the source dataset and its authors.} An annotator that labels the existing images 20 times more densely in five minutes, in the dataset's own formats, and the written operational definitions the annotation process never had. The fitted \texttt{near} threshold is a direct answer to the future-work request in Wang et al. (2025) for "spatial thresholds for near". Two annotation defects are quantified for the first time: an inverted front/behind convention in two of nine annotator groups, and \texttt{near} used by only three groups with fourfold variation in exhaustiveness.

\textbf{For work on scene-graph benchmarks.} Evidence that ranked recall against sparse, guideline-free annotation partly measures agreement with annotator habits rather than spatial correctness. The evidence is a dissociation: the model that ranks better memorises which pairs annotators record, and the model that covers the relation types they never recorded is the one trained on consistent computed labels. The per-annotator decomposition localises the effect precisely, and the seed replication shows which parts of it survive training variance.

\textbf{For weak supervision.} A controlled three-way comparison, on the same features, model, split and seeds, showing that programmatic labels out-teach both scarce human labels and the standard remedy for scarce human labels. The mechanism is measured rather than inferred: self-training contributes roughly a thousand confident negative pseudo-labels for every positive one, propagating the annotators' silence rather than their judgement.

\textbf{Methodologically.} Calibration held out by \textit{annotator} rather than by image; a sparse-gold evaluation protocol that pairs recall with audited precision; exhaustive loss attribution; a way to estimate what annotators would score against one another when they never labelled the same images, by using a deterministic annotator as a fixed common reference (\S4.6); and a reliability check that needs no labels at all, obtained by recovering the fact that an image dataset was cut from a continuous capture and asking whether its labels survive the camera moving (\S4.12). The last of these is the one most likely to transfer. Many robotics datasets are sequences presented as image sets, and wherever that is true, a predicate's agreement with itself across viewpoints separates a rule that is wrong from a rule that is merely uncertain, which is a distinction sparse human annotation cannot draw.

\section{Limitations, and future research and development}

Each limitation below is stated with the specific experiment that would settle it, because that is more useful than an apology.

\textbf{The chain reaches the plan, not the robot.} The planner experiment (\S5.7) closes one of the two remaining links: across 25 held-out scenes an LLM planner never clears an occluding object when given objects alone (0/25), always clears it when given the human relationships (25/25), and does so on 22 of 25 with the automatic ones, all three failures traced to a missing support relation rather than to faulty reasoning. What is missing is execution. No robot moved during this project, so the evidence runs from labels to models to plans and stops. Putting the same conditions on a physical Spot, or in a simulator with contact physics, would close the last link, and it is now the only one left in the chain.

\textbf{Precision estimates remain partly author-verdicted.} The independent validation study (Appendix E.3) is deployed and collecting, with coverage stratified so the audit-overlap items reach several raters first. Until it completes the audited figures carry the author's own judgement, conservatively applied and published for checking.

\textbf{One domain, one camera, six object classes.} What transfers is the calibration procedure, not the fitted constants. A labelled cross-domain sample of a few dozen images would turn Appendix E.4's qualitative evidence into a measurement, and is cheap enough that a replication should simply include one.

\textbf{Front/behind is bounded, but less by depth than the number suggests.} Section 4.9 bounds the engineering: neither a larger depth model nor multi-frame geometry moves the pair, so the limit is monocular ambiguity in the scenes rather than model capacity. What that leaves is a limitation of a different kind, since a predicate reproducing its own verdict across viewpoints 0.955 of the time while recovering 0.64 of the human labels is not mismeasuring the scene but applying a criterion the annotators did not share. The intervention with the best expected return is therefore a written annotation guideline rather than a better network, which is an uncomfortable conclusion for a computer-vision project and the one the evidence supports. The measurement routes that stay open are a calibrated stereo pair or an RGB-D capture; wider surface detection was built, measured and declined (Appendix D.4).

\textbf{Detection bounds full automation.} End-to-end recall with a zero-shot detector is 0.38 against the relation layer's 0.85 conditional on detection. The gap is detection, not relations, and the source paper's own trained detector (0.93 mAP@50) would close most of it. That check needs only the released weights.

\textbf{Scale is demonstrated without ground truth.} The supervising group supplied the full 2,650-frame capture the released images were cut from, and the pipeline was run over the 1,766 frames nobody has annotated: 562 keyframes after content-adaptive selection, 58 minutes, 185,242 triplets, a predicate distribution 0.032 in total variation from the annotated portion (Appendix E.6). Capacity and stability on unfamiliar input are therefore measured rather than argued. Correctness on that portion is not and cannot be without labels; \S4.12's viewpoint-consistency check substitutes self-agreement for truth and should be read as the weaker thing it is. Closing it needs a modest labelled sample from those frames, a few hundred triplets, which is an afternoon of annotation rather than a research programme.

\textbf{The capture is stereo, and only one eye was used.} The supplied folder is named \texttt{rightimg}, implying a left counterpart held by the supervising group. True stereo would attack the front/behind bound directly, supplying disparity at every frame from a known fixed baseline, which is exactly what the multi-frame estimators of A9 lack: those must recover the camera's motion first, and on small low-texture objects they answer for only 9\% of pairs and are 0.17 less accurate than the monocular cascade where they do (\S4.9, Appendix D.6). A calibrated stereo pair removes both problems at once and is the cheapest experiment left on this predicate. It also keeps the method's premise intact, since stereo is available at capture time, whereas depth recovered from a robot walking twenty frames is not available to a single-image annotator at all.

\section{Personal reflections}

The most useful thing this project taught me was to distrust a number until I know how it was produced.

That lesson arrived early and painfully. My first depth-based results were poor, and the unit tests all passed. The cause was that the dataset's images are stored rotated 180 degrees behind an EXIF flag, so every mask and depth value was being read from an upside-down image while the boxes stayed upright. No automated check could have caught it, because every component was behaving exactly as written. I found it by rendering an image with its boxes drawn on and looking at it. Since then I have rendered and inspected samples after every significant change, and it has caught things twice more.

The second lesson was that "agreement with the humans" is not one target. I began by treating the human labels as ground truth and my disagreements as errors. Measuring them properly showed that three of nine annotator groups used \texttt{near} at all, that two labelled front and behind in opposite directions, and that some labelled support in only one direction. That reframed the project: the interesting question stopped being "how close can I get to the humans" and became "what exactly do the humans mean, and where do they disagree with each other". Almost every later design decision, including calibrating on some annotators and validating on others, follows from that shift.

The third lesson was about my own claims, and the one I would most like to have learned sooner. On a single benchmark run I reported that my labels won against the one test annotator with clean conventions, on a margin of 0.011 over 73 images; two more seeds showed them tied and I withdrew it (\S6.3.1). I had been careful in one place, fitting thresholds on some annotators and validating on others, and careless in another, treating one training run as a result. The discipline was not new to me, I had simply not applied it to model training variance, and recording the retraction rather than quietly deleting it felt uncomfortable and is, I think, right.

Given the time again I would front-load the experiments that answer the question the project actually asks. I spent considerable effort proving the labels are correct and train models well, and comparatively little on whether they help a robot decide what to do, which is the entire motivation. The planner experiment ran last and was the clearest single result: 0 of 25 against 25 of 25 needs no statistics, and run in week five it would have pointed the intervening work at supply of the support relation, where its only failures came from. I would also have had two people annotate the same fifty images in week two, since one afternoon would have produced outright the inter-annotator agreement figure I ended up estimating indirectly.

Technically, the thing I am most pleased with is not the accuracy figure but the two refinements I built, measured and then declined. Both were plausible, both took real work, and the data said neither helped. Reporting them as null results rather than dropping them was the point at which the project started feeling like research rather than engineering.